\section{Autoregressive models strike back}

\subsection{为什么 AR 又回来了}

raw pixels 上的 autoregressive 模型太慢了，因为序列太长。但如果把图像先压到离散 latent 空间，AR 就重新变得可行。也就是说，AR 不是被 diffusion 完全淘汰了，而是换了赛道：它更适合做离散 token 的序列建模。

\begin{figure}[H]
\centering
\includegraphics[width=0.95\textwidth]{slides-images/slide-119.jpg}
\caption{Autoregressive models strike back：它们在 raw pixels 上太慢，但在离散 latent 上又重新变得好用\protect\footnotemark}
\end{figure}
\footnotetext{Slides 第120页。}

\begin{figure}[H]
\centering
\includegraphics[width=0.95\textwidth]{slides-images/slide-120.jpg}
\caption{离散 latent + AR：先学 encoder/decoder，再让 autoregressive 模型预测离散码本索引\protect\footnotemark}
\end{figure}
\footnotetext{Slides 第121页。}

\begin{importantbox}{为什么离散 latent 很有吸引力}
离散 latent 把图像压缩成更短的 token 序列，序列建模任务就从“生成几十万像素”变成“生成一小串离散码”。这对 Transformer 特别友好，也让高质量文本到图像生成再次有了 AR 路线。
\end{importantbox}

